% Source: https://artofproblemsolving.com/community/c3h2527370_stuffedmaths_problem_thread
% Stable identity: aops-topic-2527370-q22
\emph{Top and bottom.} Seen from above, the exposed horizontal faces exactly cover the $6\times6$ top of the bottom cube (each cube covers part of the top of the cube below, and its own top is exposed). So the upward-facing area is $36$, and the bottom face adds another $36$: total $72$.

\emph{Sides.} No cube's vertical faces are covered: a cube sits on top of a larger one and only touches it along its bottom. So every cube contributes all four of its side faces:
\[4\left(1^2+2^2+3^2+4^2+5^2+6^2\right)=4\cdot91=364.\]
(Faces that are flush with faces of the cube below are still exposed; they just lie in the same plane.)

The surface area is $72+364=\boxed{436}$.

\begin{asy}
usepackage("amsmath");
size(4.5cm);
defaultpen(fontsize(10pt));
pair P(real x,real y,real z){return (x+0.45*y,z+0.32*y);}
void cube(real s, real z0){
  real x0=6-s, y0=0;
  path f1=P(x0,y0,z0)--P(6,y0,z0)--P(6,y0,z0+s)--P(x0,y0,z0+s)--cycle;
  path f2=P(6,y0,z0)--P(6,s,z0)--P(6,s,z0+s)--P(6,y0,z0+s)--cycle;
  path f3=P(x0,y0,z0+s)--P(6,y0,z0+s)--P(6,s,z0+s)--P(x0,s,z0+s)--cycle;
  filldraw(f1,white,linewidth(0.7));filldraw(f2,lightgray,linewidth(0.7));filldraw(f3,gray(0.9),linewidth(0.7));
}
real z=0;
for(int s=6;s>=1;--s){cube(s,z); z+=s;}
label("$6$",P(3,0,0),S);label("$6$",P(6,3,0),SE);
\end{asy}
